(1) 在指定规范化Fock基中 \(H=\begin{pmatrix}U&-\sqrt2J&0\\-\sqrt2J&0&-\sqrt2J\\0&-\sqrt2J&U\end{pmatrix}\)。因 \(a_L^\dagger a_R|1,1\rangle=\sqrt2|2,0\rangle\)，非对角元必须带 \(\sqrt2\)，不能当成一粒子隧穿系数 \(-J\)。

(2) 交换反对称态 \(|A\rangle=(|2,0\rangle-|0,2\rangle)/\sqrt2\) 与 \(|1,1\rangle\) 完全解耦，能量为 \(U\)。对称态 \(|S\rangle=(|2,0\rangle+|0,2\rangle)/\sqrt2\) 与 \(|1,1\rangle\) 的块为 \(\begin{pmatrix}U&-2J\\-2J&0\end{pmatrix}\)，能级 \(\boxed{E_\pm=(U\pm\sqrt{U^2+16J^2})/2}\)，基态为 \(E_-<0\)。

(3) 令 \(\Delta=\sqrt{U^2+16J^2}\)，Hellmann–Feynman对 \(U\) 求导得到基态双占据概率 \(P_d=\langle[n_L(n_L-1)+n_R(n_R-1)]/2\rangle=\partial_UE_-=(1-U/\Delta)/2\)。对称性保证 \(\langle n_L\rangle=1\)，且 \(\operatorname{Var}(n_L)=P_d\)。\(U=0\) 时 \(P_d=1/2\)，强排斥 \(U\gg J\) 时 \(P_d\simeq4J^2/U^2\)。

(4) 对 \(J\) 求导 \(\partial_JE_-=-8J/\Delta=-\langle a_L^\dagger a_R+a_R^\dagger a_L\rangle\)，在实对称基态下得到 \(\boxed{\langle a_L^\dagger a_R\rangle=4J/\Delta}\)。无相互作用时为1，强排斥时约 \(4J/U\) 趋零，反映数目固定与位点相干性之间的量子竞争。
